\expandafter\def\csname HistoricalNoteText404\endcsname{Resultate, 1837. IV. [S. 383. d. B.]}
\renewcommand{\Article}[1]{\par\medskip\Needspace{4\baselineskip}\pdfbookmark[1]{#1.}{ablenkungen-art-#1}\centerline{\textbf{#1.}}\nobreak\smallskip}%
\SourceLeaf{402}werden in solchen Fällen, wo die Stärke der Zusatzkraft schon in kurzer Zeit beträchtliche Veränderungen erleidet, wie bei galvanischen Strömen.

\Article{2}
Durch folgendes einfache Verfahren wird diesem Übelstande abgeholfen. Man lässt die Zusatzkraft zuerst nur während des dritten Theils der Schwingungsdauer wirken, suspendirt sie dann während einer eben so langen Zwischenzeit, und setzt sie darauf erst in beharrliche Wirksamkeit. Ist also z. B. die Schwingungsdauer der Nadel des Magnetometers 30 Secunden, und soll die durch einen galvanischen Strom erzeugte Ablenkung gemessen werden, so schliesst man die Kette bei einem Secundenschlage, welchen man als 0 zählt; öffnet wieder bei $10''$, und schliesst endlich definitiv bei $20''$. Soll die Ablenkung durch einen an einen bestimmten Platz zu legenden Magnetstab geschehen, so nähert man sich dem vorher genau und bequem bezeichneten Platze mit dem anfangs vertical gehaltenen Magnetstabe, legt denselben bei $0''$ plötzlich nieder, richtet ihn bei $10''$ eben so schnell wieder auf und legt ihn zum zweiten Male bei $20''$ definitiv hin. Der Erfolg ist, dass die Nadel von ihrer ursprünglichen Ruhestellung sich derjenigen Stellung, welche der Ablenkung entspricht, während der ersten 10 Secunden mit beschleunigter Geschwindigkeit nähert, bei $10''$ gerade die Mitte zwischen beiden Stellungen erreicht hat, und dann während der zweiten 10 Secunden die andere Hälfte des Zwischenraumes mit retardirter Bewegung durchläuft, so dass sie bei 20 Secunden die neue Stellung erreicht und alle Bewegung verloren haben wird.

Man sieht leicht, dass auf ganz ähnliche Weise die Nadel von einem ruhigen Ablenkungszustande zu dem entgegengesetzten so hinübergeführt werden kann, dass sie in demselben ohne Bewegung ist: man lässt nemlich die ablenkende Kraft während des dritten Theils der Schwingungsdauer im entgegengesetzten Sinne wirken, dann während eben so langer Zeit wieder im frühern Sinn, und wechselt darauf von neuem. Für galvanische Ströme erhält man den Wechsel fast augenblicklich durch einen zweckmässigen Commutator; für ablenkende Magnetstäbe durch eine rasche halbe Umdrehung (am bequemsten durch eine horizontale), so dass der Nordpol des Stabes an den Platz des Südpols kommt.

Endlich ist klar, dass auf dieselbe Weise nach beobachteter Ablenkung die Nadel wieder ruhig in den reinen magnetischen Meridian gebracht werden kann: man braucht nur die ablenkende Kraft zuerst während eines Drittheils der Schwin\SourceLeaf{403}gungsdauer zu suspendiren, dann eben so lange noch einmal wirken und endlich aufhören zu lassen.

\Article{3}
Dem beschriebenen Verfahren liegt die Voraussetzung zum Grunde dass

\textit{erstens} die Schwingungen der Nadel so erfolgen, dass der Abstand von der Mitte der Schwingung (so lange diese Mitte selbst nicht abgeändert wird) dem Sinus eines sich gleichförmig ändernden und während einer Schwingungsdauer um $180^0$ zunehmenden Winkels proportional bleibt, und

\textit{zweitens}, die Schwingungsdauer durch die Zusatzkraft nicht verändert wird.

Insofern beide Voraussetzungen nicht \textit{in absoluter Schärfe} gültig sind, und ausserdem auch bei der Ausführung weder der Wechsel ganz augenblicklich geschehen, noch die vorgeschriebenen Zwischenzeiten absolut genau eingehalten werden können, wird allerdings nach Vollendung der Operation die Nadel selten in vollkommener Ruhe angetroffen werden: allein für den praktischen Zweck ist es schon hinreichend, wenn die übrig bleibende Bewegung so gering ist, dass man die wahre Gleichgewichtsstellung auf gewöhnliche Weise sogleich zu beobachten anfangen kann.

Unter den Statt findenden Umständen werden jene Voraussetzungen nur sehr wenig von der Wahrheit abweichen können. Die Anwendbarkeit des Magnetometers beruht an sich schon darauf, dass die Zusatzkraft nur eine mässige Ablenkung hervorbringt, wobei (einen sogleich zu erwähnenden Ausnahmefall beiseite gesetzt) das in der ersten Voraussetzung enthaltene Gesetz hinreichend genau gilt. Die Veränderung der Schwingungsdauer durch die ablenkende Kraft ist ganz unmerklich, wenn diese senkrecht gegen den magnetischen Meridian wirkt, wie fast immer der Fall ist: wirkte sie aber auch in einer schiefen Richtung, so würde, insofern sie selbst nur ein kleiner Bruchtheil der erdmagnetischen Kraft ist, die dadurch bewirkte Veränderung der Schwingungsdauer doch für die kurze Zeit der Operation ganz unerheblich bleiben.

Nur Ein Fall ist auszunehmen, nemlich wenn die Nadel ihre Schwingungen unter dem Einflusse eines die Grösse des Schwingungsbogens bedeutend vermindernden Dämpfers macht. In diesem Falle ist das obige Gesetz nicht mehr gültig, und eine genaue Befolgung des oben beschriebenen Verfahrens würde nicht zum Ziele führen: von der andern Seite ist dann aber auch allerdings der
\SourceLeaf{404}im 1. Art. bemerkte Übelstand viel geringer, da ein kräftiger Dämpfer die Nadel von selbst in mässiger Zeit zur Ruhe bringt. Da indessen für diesen Fall jenes Verfahren nur einer Modification bedarf, um denselben Erfolg zu erreichen, und es allemal erwünscht sein muss, jeden unnöthigen Zeitverlust vermeiden zu können, so ist es, in praktischer wie in theoretischer Beziehung, der Mühe werth. die Frage ganz allgemein zu betrachten.

\Article{4}
Wir haben zuvörderst folgende allgemeine Aufgabe aufzulösen.

Ein Magnetstab schwingt unter wiederholter Abänderung der auf ihn wirkenden Kräfte, wobei jedoch die Schwingungsdauer und das logarithmische Decrement\HistoricalNote{404} unverändert, und die Schwingungsbogen klein genug bleiben, um Grössen der dritten Ordnung vernachlässigen zu können. Man soll aus dem anfänglichen Bewegungszustande denjenigen, welcher nach der letzten Abänderung Statt findet, ableiten.

Es sei $T$ die Schwingungsdauer, $\epsilon$ das logarithmische Decrement, $e$ die Basis der hyperbolischen, $m$ der Modulus der briggischen Logarithmen, $\pi$ das Verhältniss des Kreisumfanges zum Durchmesser. Man setze
\[
n=\frac{\pi}{T},\qquad \epsilon=\frac{\lambda}{mT}
\]
Unter obigen Voraussetzungen wird demnach der Stand $x$ für die Zeit $t$ durch die Formel ausgedrückt
\[
x=p+Ae^{-\epsilon t}\sin(nt-B)
\]
welcher man auch die Gestalt geben kann
\[
x=p+ae^{-\epsilon t}\cos nt+be^{-\epsilon t}\sin nt
\]
wo $p$ die Gleichgewichtsstellung ausdrückt, und die Coëfficienten $a,b$ so lange constant bleiben, als $p$ constant ist. Die Geschwindigkeit der Bewegung findet sich hieraus
\[
\frac{dx}{dt}=-e^{-\epsilon t}(na\sin nt+\epsilon a\cos nt-nb\cos nt+\epsilon b\sin nt)
\]
oder wenn man einen Hülfswinkel $\varphi$ einführt, so dass $\frac{\epsilon}{n}=\operatorname{tang}\varphi$ wird,
\SourceLeaf{405}\[
\frac{dx}{dt}=-\frac{ne^{-\epsilon t}}{\cos\varphi}(a\sin(nt+\varphi)-b\cos(nt+\varphi))
\]
Für $ae^{-\epsilon t}\cos nt+be^{-\epsilon t}\sin nt$ schreiben wir $u$, so dass $x=p+u$ wird.

Es seien nun $t',t'',t'''$ die bestimmten Werthe von $t$, wo eine Veränderung der wirkenden Kraft vorgenommen wird; ferner seien die bestimmten Werthe von $p,a,b$ in den verschiedenen Zeitabschnitten folgende:
\[
\begin{array}{lll}
p^0,\ a^0,\ b^0&\text{vor}&t'\\
p',\ a',\ b'&\text{von}&t'\ \text{bis}\ t''\\
p'',\ a'',\ b''&\text{von}&t''\ \text{bis}\ t'''\\
p''',\ a''',\ b'''&\text{nach}&t'''
\end{array}
\]
Endlich gehe der allgemeine Ausdruck von $u$, wenn für $a$ und $b$ die bestimmten Werthe substituirt werden, über in $u^0,u',u'',u'''$, so dass vor dem ersten Wechsel $x=p^0+u^0$ wird, von da bis zum zweiten $x=p'+u'$ u. s. f.

Da der Augenblick $t'$ zugleich der letzte des ersten Zeitabschnitts und der erste des folgenden ist, so müssen für $t=t'$ sowohl $x$ als $\frac{dx}{dt}$ einerlei Werth erhalten, man möge in den obigen allgemeinen Ausdrücken für $p,a,b$ die Werthe $p^0,a^0,b^0$, oder $p',a',b'$ substituiren.

Es ist also
\[
\begin{aligned}
0&=p'-p^0+(a'-a^0)e^{-\epsilon t'}\cos nt'+(b'-b^0)e^{-\epsilon t'}\sin nt'\\
0&=(a'-a^0)\sin(nt'+\varphi)-(b'-b^0)\cos(nt'+\varphi)
\end{aligned}
\]
woraus man leicht ableitet
\[
\begin{aligned}
a'-a^0&=-\frac{p'-p^0}{\cos\varphi}.e^{\epsilon t'}\cos(nt'+\varphi)\\
b'-b^0&=-\frac{p'-p^0}{\cos\varphi}.e^{\epsilon t'}\sin(nt'+\varphi)
\end{aligned}
\]
und hieraus
\[
u'=u^0-\frac{p'-p^0}{\cos\varphi}.e^{-\epsilon(t-t')}\cos(n(t-t')-\varphi)
\]
Auf gleiche Art erhält man
\[
\begin{aligned}
u''&=u'-\frac{p''-p'}{\cos\varphi}.e^{-\epsilon(t-t'')}\cos(n(t-t'')-\varphi)\\
u'''&=u''-\frac{p'''-p''}{\cos\varphi}.e^{-\epsilon(t-t''')}\cos(n(t-t''')-\varphi)
\end{aligned}
\]
und so ferner, wenn noch mehrere Wechsel der bewegenden Kräfte Statt finden.

\SourceLeaf{406}Es wird also hiedurch aus dem anfänglichen Bewegungszustande jeder nachfolgende bestimmt.

\Article{5}
Für den Fall der gegenwärtigen Untersuchung ist $p''=p^0$ und $p'''=p'$ zu setzen. Dadurch wird
\[
\begin{aligned}
u'''={}&u^0-\frac{p'-p^0}{\cos\varphi}.e^{-\epsilon t}.\\
&.\{e^{\epsilon t'}\cos(n(t-t')-\varphi)-e^{\epsilon t''}\cos(n(t-t'')-\varphi)\\
&\qquad+e^{\epsilon t'''}\cos(n(t-t''')-\varphi)\}
\end{aligned}
\]
welche Formel, wenn man
\[
\begin{aligned}
e^{-\epsilon(t''-t')}\cos n(t''-t')-1+e^{\epsilon(t'''-t'')}\cos n(t'''-t'')&=f\\
e^{-\epsilon(t''-t')}\sin n(t''-t')\phantom{-1}{}-e^{\epsilon(t'''-t'')}\sin n(t'''-t'')&=g
\end{aligned}
\]
setzt, übergeht in
\[
\begin{aligned}
u'''={}&u^0-\frac{p'-p^0}{\cos\varphi}.e^{-\epsilon(t-t'')}.\{f\cos(n(t-t'')-\varphi)\\
&\qquad-g\sin(n(t-t'')-\varphi)\}
\end{aligned}
\]
Hieraus folgt, dass wenn die Zwischenzeiten $t''-t',t'''-t''$ so bestimmt sind, dass $f=0$ und $g=0$ wird, allgemein
\[
u'''=u^0
\]
oder $a'''=a^0,b'''=b^0$ wird.

War also vor den Wechseln die Nadel in $p^0$ in Ruhe, so wird sie, nach denselben, sich in $p'$ in Ruhe befinden: im entgegengesetzten Falle wird die Nadel nach den drei Wechseln in jedem Augenblick genau dieselbe Geschwindigkeit und dieselbe Stellung gegen den Mittelpunkt ihrer Bewegung $p'$ haben, welche sie relativ gegen $p^0$ in demselben Augenblicke haben würde, wenn sie ihre ursprüngliche Bewegung ungestört fortgesetzt hätte: mit Einem Worte, bloss der Mittelpunkt der Bewegung wird versetzt, die Bewegung selbst aber gar nicht geändert sein.

\Article{6}
Es bleibt nun noch übrig, die Zwischenzeit so zu bestimmen, dass den Gleichungen $f=0,g=0$ Genüge geschehe. Setzt man
\SourceLeaf{407}\[
t''-t'=qT,\qquad t'''-t''=rT
\]
und erinnert sich, dass $e=10^m$, so werden jene Gleichungen
\[
\begin{aligned}
10^{-q\lambda}\cos q\pi+10^{r\lambda}\cos r\pi&=1\\
10^{-q\lambda}\sin q\pi&=10^{r\lambda}\sin r\pi
\end{aligned}
\]
Für den Fall einer unmerklichen Abnahme des Schwingungsbogens muss also $\cos q\pi+\cos r\pi=1$ und $\sin q\pi=\sin r\pi$ gesetzt werden, mithin
\[
q\pi=r\pi=60^0\ \text{oder}\ =\tfrac{1}{3}\pi,\quad\text{und}\quad t''-t'=t'''-t''=\tfrac{1}{3}T
\]
wie schon im 2. Artikel bemerkt ist. Für den Fall eines merklichen logarithmischen Decrements hingegen werden jene Gleichungen auf indirectem Wege aufzulösen sein, welcher Rechnung man folgende Form geben kann.

Aus der Verbindung der Gleichungen folgt
\[
\begin{gathered}
\operatorname{tang}r\pi=\frac{\sin q\pi}{10^{q\lambda}-\cos q\pi},\qquad
10^{2r\lambda}=1-2.10^{-q\lambda}\cos q\pi+10^{-2q\lambda}
\end{gathered}
\]
Durch Elimination von $r$ hat man also die Gleichung mit Einer unbekannten Grösse
\[
\frac{\pi}{2\lambda}\log(1-2.10^{-q\lambda}\cos q\pi+10^{-2q\lambda})=\operatorname{Arc.\,tang}\frac{\sin q\pi}{10^{q\lambda}-\cos q\pi}
\]
wo der briggische Logarithme verstanden ist. Nachdem derselben Genüge geleistet ist, hat man offenbar zugleich den Werth von $r$.

\Article{7}
Um denjenigen, welche das beschriebene Verfahren unter Anwendung eines Dämpfers ausüben wollen, die im vorhergehenden Artikel erklärte Rechnung zu ersparen, theile ich hier eine im voraus berechnete Tafel mit, aus welcher für jedes logarithmische Decrement das Verhältniss der beiden Zwischenzeiten zur Schwingungsdauer sogleich entnommen werden kann. Man sieht daraus, dass mit zunehmendem logarithmischen Decrement die erste Zwischenzeit immer grösser, die zweite immer kleiner wird. Die Summe beider ist zwar zwei Drittheilen der Schwingungsdauer nur für $\lambda=0$ genau gleich, entfernt sich aber davon viel langsamer. Dass es bei der wirklichen Anwendung zureicht, etwa nur die ersten Decimalen der Werthe von $q$ und $r$ zu berücksichtigen, bedarf keiner Erinnerung.

\SourceLeaf{408}\begin{center}\textit{\textbf{Tafel.}}\end{center}
\begingroup\fontsize{9}{11}\selectfont
\begin{center}
\begin{tabular}{c|c|c||c|c|c}
$\lambda$&$q$&$r$&$\lambda$&$q$&$r$\\\hline
0&0.33333&0.33333&0.30&0.45921&0.21406\\
0.01&0.33757&0.32911&0.31&0.46322&0.21048\\
0.02&0.34181&0.32489&0.32&0.46721&0.20694\\
0.03&0.34606&0.32068&0.33&0.47118&0.20343\\
0.04&0.35031&0.31648&0.34&0.47513&0.19996\\
0.05&0.35456&0.31229&0.35&0.47906&0.19652\\[3pt]
0.06&0.35882&0.30812&0.36&0.48297&0.19311\\
0.07&0.36308&0.30395&0.37&0.48685&0.18975\\
0.08&0.36734&0.29981&0.38&0.49071&0.18641\\
0.09&0.37160&0.29568&0.39&0.49454&0.18311\\
0.10&0.37585&0.29156&0.40&0.49835&0.17985\\[3pt]
0.11&0.38011&0.28746&0.41&0.50214&0.17663\\
0.12&0.38436&0.28338&0.42&0.50590&0.17344\\
0.13&0.38861&0.27932&0.43&0.50963&0.17029\\
0.14&0.39285&0.27528&0.44&0.51334&0.16718\\
0.15&0.39708&0.27126&0.45&0.51702&0.16411\\[3pt]
0.16&0.40131&0.26727&0.46&0.52067&0.16107\\
0.17&0.40552&0.26329&0.47&0.52430&0.15808\\
0.18&0.40973&0.25934&0.48&0.52790&0.15512\\
0.19&0.41393&0.25542&0.49&0.53147&0.15220\\
0.20&0.41812&0.25152&0.50&0.53501&0.14931\\[3pt]
0.21&0.42230&0.24764&0.51&0.53852&0.14647\\
0.22&0.42646&0.24379&0.52&0.54201&0.14367\\
0.23&0.43061&0.23997&0.53&0.54546&0.14090\\
0.24&0.43474&0.23618&0.54&0.54889&0.13817\\
0.25&0.43886&0.23242&0.55&0.55229&0.13548\\[3pt]
0.26&0.44297&0.22868&0.56&0.55566&0.13283\\
0.27&0.44705&0.22498&0.57&0.55900&0.13022\\
0.28&0.45112&0.22131&0.58&0.56231&0.12765\\
0.29&0.45517&0.21767&0.59&0.56559&0.12511\\
0.30&0.45921&0.21406&0.60&0.56884&0.12261
\end{tabular}
\end{center}\endgroup

\Article{8}
Die unserer Theorie zum Grunde liegende Voraussetzung, dass die drei Wechsel augenblicklich geschehen, findet bei der wirklichen Ausübung des Verfahrens in aller Schärfe niemals Statt, obwohl bei Ablenkungen durch galvanische Ströme die zu jedem Wechsel nöthige Zeit als unmerklich betrachtet werden kann. Bei Ablenkungen durch Magnetstäbe hingegen ist, nach Maassgabe ihrer Grösse und Schwere, diese Zeit schon mehr oder weniger bedeutend, und bei fünfundzwanzigpfündigen werden zu Vollführung eines Wechsels immer mehrere Secunden erforderlich sein, besonders wenn nicht von einem Wechsel zwischen verticaler und horizontaler Lage, sondern zwischen zweien entgegengesetzten La\SourceLeaf{409}gen die Rede ist. Für diesen Fall, welcher in der That der bei weiten wichtigste und gewöhnlichste ist, lässt sich aber die Ausführung der Operation leicht so einrichten, dass der Erfolg kaum merklich gestört wird. Man muss nur Sorge tragen, dass der zweite und dritte Wechsel auf gleiche Weise geschehen, wie der erste, also auch eine gleich lange Zeit ausfüllen, und diese Zeit den sonst nöthigen Zwischenzeiten abbrechen. Ist z. B. (wie \textit{Res.} 1837. IV.) [S. 390. d. B.] das logarithmische Decrement 0.33570, die Schwingungsdauer $21''21439$, so folgt aus obiger Tafel die erste Zwischenzeit $=10''04$, die zweite $=4''27$: findet man nun zur Ausführung eines Wechsels drei Secunden nöthig, so \textit{beginnt} man den ersten Wechsel bei $0''$; von $3''$ bis $10''$. bleibt der Stab in der neuen Lage; durch den bei $10''$ \textit{anfangenden} neuen Wechsel ist der Stab bei $13''$ in die erste Lage zurückgebracht, in welcher er nur $1\tfrac{1}{4}$ Secunden liegen bleibt, worauf der dritte Wechsel anfängt, so dass erst mit $17\tfrac{1}{4}$ Secunden die ganze Operation vollendet ist. Eine ausgedehntere, hier jedoch des Raumes wegen zu übergehende Untersuchung ergibt nemlich, dass wenn $p^0$ in $p'$ nicht sprungsweise sondern allmählig übergeht, und eben so beim zweiten Wechsel $p'$ in $p^0$, und beim dritten wiederum $p^0$ in $p'$, der Erfolg ganz derselbe bleibt, wie er am Schluss des 5. Artikels angegeben ist, falls nur die drei Übergangszeiten gleich lang sind, die drei Übergänge selbst in ähnlichen Stufenfolgen geschehen, und die berechneten Zwischenzeiten $qT,rT$ auf die Anfangsmomente der Wechsel bezogen, oder was dasselbe ist die beiden ersten Übergangszeiten ihnen eingerechnet werden.

\clearpage
\fancyhead[L]{\small CONSTANTEN DES BIFILARMAGNETOMETERS.}
\SourceLeaf{410}\pdfbookmark[0]{Zur Bestimmung der Constanten des Bifilarmagnetometers.}{bifilarconstanten-title}
\begin{center}
{\Large\bfseries ZUR BESTIMMUNG\par\medskip DER CONSTANTEN DES BIFILARMAGNETOMETERS.}
\par\bigskip
Resultate aus den Beobachtungen des magnetischen Vereins. 1840. I.
\end{center}
\renewcommand{\Article}[1]{\par\medskip\Needspace{4\baselineskip}\pdfbookmark[1]{#1.}{bifilarconstanten-art-#1}\centerline{\textbf{#1.}}\nobreak\smallskip}%
\Article{1}
Zum richtigen und sichern Gebrauche des Bifilarmagnetometers ist die Kenntniss der Zahlenwerthe gewisser Grössen erforderlich, die sich auf bedingungsweise wie constant zu betrachtende Verhältnisse der Theile des Apparats beziehen, und von denen als wesentlichen Elementen die nach den verschiedenen Stellungen der beweglichen Theile zu beobachtenden Gleichgewichtslagen und Schwingungszeiten abhängen. Diese Elemente sind vier, nemlich

1) die Stellung, welche der Index der Spiegelalhidade haben muss, damit die Normale gegen den Spiegel mit der optischen Axe des Beobachtungsfernrohrs in Eine Verticalebene falle, wenn die beiden Aufhängungsdrähte in einer Verticalebene sind; diese Stellung (so verstanden, dass die reflectirende Fläche des Spiegels dem Fernrohre zugekehrt sei) soll mit $P$ bezeichnet werden.

2) die Stellung, welche bei eben dieser Lage der Aufhängungsdrähte dem Index des Schiffchens gegeben werden muss, damit die magnetische Axe des Magnetstabes sich in natürlicher Lage im magnetischen Meridiane befinde; ich bezeichne diese Stellung mit $Q$.

Es bedarf keiner Erinnerung, dass wenn jede der beiden Alhidaden mehr als einen Index hat, einer davon immer (nach Belieben) als Hauptindex zu wählen ist.

3) das Verhältniss der magnetischen Directionskraft zu der aus der Aufhängungsweise entspringenden, welche letztere die statische Directionskraft heissen mag: dieses Verhältniss soll durch $R:1$ ausgedrückt werden.

\SourceLeaf{411}4) die statische Schwingungsdauer des Apparats, d. i. diejenige, welche bloss in Folge der Aufhängungsart oder ohne Einwirkung des Erdmagnetismus auf den Magnetstab, Statt finden würde: ich bezeichne das Quadrat dieser Schwingungsdauer mit $S$.

Es erhellt hieraus, dass $\frac{S}{R}$ das Quadrat der reinmagnetischen Schwingungsdauer ausdrückt, d. i. derjenigen, die bei der Aufhängung des Apparats an einem einfachen Faden ohne Torsion Statt haben würde.

\Article{2}
Es ist nun zuvörderst zu entwickeln, wie das, was am Bifilarmagnetometer unmittelbar beobachtet wird, mit der Stellung der beiden Alhidaden und diesen vier Elementen zusammenhängt.

Bei der Stellung der Alhidade des Spiegels auf $A$, der Alhidade des Schiffchens auf $B$, bezeichne $t$ die Schwingungsdauer, und $p$ den in Bogentheile verwandelten Abstand des der Gleichgewichtslage entsprechenden Skalentheils von demjenigen Punkte der Skale, der mit der optischen Axe des Beobachtungsfernrohrs in derselben Verticalebene ist, und durch den von der Mitte des Objectivs herabhängenden Lothfaden kenntlich gemacht wird. Um die Vorstellungen zu fixiren, nehme ich an, dass die Theilungen sowohl am Kreise als an der Skale von der Linken nach der Rechten laufen, und beziehe positive Zeichen von $p$ auf den Fall, wo die auf dem Fadenkreuze des Fernrohrs beobachtete Zahl grösser ist, als die Zahl am Lothfaden. Bei jener Gleichgewichtslage befindet sich also das Bifilarmagnetometer um $A-P-p$ rückwärts, d. i. von der Rechten nach der Linken gedreht gegen diejenige Lage, wo die Aufhängungsdrähte parallel waren, oder $A-P-p$ ist der Winkel zwischen der geraden Linie durch die beiden untern Enden der Aufhängungsdrähte und einer Parallele mit der die beiden obern Enden verbindenden. Das durch die Aufhängungsweise hervorgebrachte Drehungsmoment ist zwar nicht in völliger Schärfe, aber hinlänglich genau für die Ausübung, dem Sinus dieses Winkels proportional; wir setzen dasselbe $=D\sin(A-P-p)$, wo also $D$ die statische Directionskraft ausdrückt: die positiven Werthe des Drehungsmoments beziehen sich auf Drehung von der Linken nach der Rechten.

In derselben Lage des Apparats macht die magnetische Axe des Magnetstabes mit dem magnetischen Meridiane den von der Rechten nach der Linken ge%
